%% Supplementary Information for the PeerJ Sports Medicine and Rehabilitation submission.
%% Standalone document, builds independently of manuscript.tex.

\documentclass[fleqn,10pt,lineno]{wlpeerj}

\usepackage{bm}
\usepackage{array}
\usepackage{tabularx}
\usepackage{textcomp}
\usepackage{multirow}
\usepackage{needspace}
\usepackage{setspace}
\doublespacing

%% Keep tables and their captions single-spaced in the review copy.
\DeclareCaptionFont{singlespaced}{\setstretch{1}}
\captionsetup{font=singlespaced}
\AtBeginEnvironment{table}{\singlespacing}

\usepackage[colorlinks=true,linkcolor=black,citecolor=black,urlcolor=blue]{hyperref}

\setlength{\parskip}{0.6em plus 0.1em minus 0.1em}
\setlength{\parindent}{0pt}

\newcommand{\dps}{$^{\circ}$/s}

% Supplementary needs no abstract: render title + authors only.
\makeatletter
\renewcommand{\@maketitle}{%
  \thispagestyle{empty}%
  \vskip-36pt%
  {\raggedright\sffamily\bfseries\fontsize{20}{25}\selectfont \@title\par}%
  \vskip10pt%
  {\raggedright\sffamily\fontsize{12}{16}\selectfont \@author\par}%
  \vskip25pt%
}
\makeatother

\title{Supplementary Information: Criterion validity of heel-mounted inertial sensors for stride-length estimation during single- and dual-task walking in older adults with and without cognitive impairment}

\author[1,*]{Pawel Kudzia}
\author[1]{Markus von Hacht}
\author[1]{Maryam Sheikhi}
\author[1]{Drew Commandeur}
\author[1,2]{Marc Klimstra}
\author[1,3]{Sandra Hundza}
\affil[1]{Motion, Mobility and Rehabilitation Lab, School of Exercise Science, Physical and Health Education, University of Victoria, Victoria, BC, Canada}
\affil[2]{Canadian Sport Institute-Pacific, Victoria, BC, Canada}
\affil[3]{Institute on Aging and Lifelong Health, Victoria, BC}
\affil[*]{Corresponding author: pawel.kudzia5@gmail.com}

\begin{document}

\flushbottom
\maketitle
\thispagestyle{empty}

\renewcommand{\thetable}{S\arabic{table}}
\renewcommand{\thefigure}{S\arabic{figure}}
\setcounter{table}{0}
\setcounter{figure}{0}

This document contains the supplementary tables referenced in the main manuscript: per-participant demographics (S1), reported-condition data availability (S2), a MoCA-based cognitive-status sensitivity reanalysis (S3), a same-pair gravity-residual correction ablation (S4), a participant-clustered bootstrap of cohort-level accuracy (S5), a sensitivity analysis of the stride-pair mismatch filter (S6), and a correction-gain sensitivity analysis (S7).

% ============================================================================
% S1: Per-participant demographics
% ============================================================================

\section*{S1. Per-participant demographics}

Table~\ref{tab:s_demographics} lists demographic characteristics for all 33 enrolled participants. We excluded one participant (ID\_191) because the available recordings used shank and hip sensors rather than the heel-mounted sensors required for this validation. The remaining 32 participants contributed every reported condition for which the paired source data were available.

\begin{table}[!ht]
\centering
\caption{Individual participant demographics. ID = numeric study code. MoCA = Montreal Cognitive Assessment total (/30). Cognitive = MoCA-defined screening group (CN = cognitively normal, MoCA $\geq$ 26; CI = cognitive impairment, MoCA 14--25). These labels do not denote clinical diagnoses. Excluded = $\times$ when the required heel-mounted recordings were unavailable.}
\label{tab:s_demographics}
\footnotesize
\begin{tabular}{@{}lccrrrcc@{}}
\toprule
ID & Age & Sex & Height (cm) & Mass (kg) & MoCA & Cognitive & Excluded \\
\midrule
ID\_194 & 60 & F & 166.0 & 89.5 & 27 & CN & \\
ID\_104 & 80 & F & 159.0 & 75.0 & 17 & CI & \\
ID\_171 & 66 & F & 168.0 & 65.4 & 30 & CN & \\
ID\_107 & 75 & F & 163.0 & 66.9 & 27 & CN & \\
ID\_196 & 71 & F & 168.4 & 70.5 & 28 & CN & \\
ID\_139 & 73 & F & 165.0 & 99.8 & 28 & CN & \\
ID\_123 & 66 & M & 177.0 & 77.6 & 25 & CI & \\
ID\_174 & 68 & M & 177.0 & 80.5 & 23 & CI & \\
ID\_183 & 60 & F & 158.5 & 67.5 & 29 & CN & \\
ID\_177 & 74 & M & 168.0 & 78.6 & 20 & CI & \\
ID\_110 & 67 & F & 163.0 & 62.6 & 30 & CN & \\
ID\_180 & 78 & M & 181.5 & 96.5 & 25 & CI & \\
ID\_185 & 67 & F & 170.0 & 61.6 & 26 & CN & \\
ID\_192 & 60 & F & 162.0 & 72.7 & 26 & CN & \\
ID\_188 & 80 & F & 160.5 & 53.1 & 21 & CI & \\
ID\_197 & 63 & F & 171.0 & 72.0 & 27 & CN & \\
ID\_189 & 63 & F & 167.5 & 83.0 & 24 & CI & \\
ID\_191 & 60 & F & 165.0 & 74.2 & 29 & CN & $\times$ \\
ID\_190 & 75 & F & 158.0 & 66.7 & 26 & CN & \\
ID\_182 & 60 & F & 173.0 & 81.7 & 28 & CN & \\
ID\_169 & 71 & M & 173.0 & 81.3 & 27 & CN & \\
ID\_186 & 66 & F & 165.0 & 61.3 & 23 & CI & \\
ID\_109 & 82 & F & 166.0 & 75.2 & 28 & CN & \\
ID\_101 & 83 & F & 158.0 & 69.6 & 28 & CN & \\
ID\_184 & 70 & F & 170.0 & 77.5 & 25 & CI & \\
ID\_126 & 83 & F & 161.0 & 50.0 & 29 & CN & \\
ID\_149 & 69 & M & 173.0 & 68.4 & 22 & CI & \\
ID\_114 & 82 & F & 163.0 & 50.4 & 26 & CN & \\
ID\_135 & 74 & F & 160.0 & 74.1 & 29 & CN & \\
ID\_193 & 75 & M & 175.0 & 79.8 & 24 & CI & \\
ID\_195 & 61 & M & 192.0 & 104.5 & 23 & CI & \\
ID\_136 & 82 & F & 161.0 & 67.4 & 28 & CN & \\
ID\_172 & 61 & F & 164.0 & 53.9 & 23 & CI & \\
\bottomrule
\end{tabular}
\end{table}

% ============================================================================
% S2: Reported-condition data availability
% ============================================================================

\section*{S2. Reported-condition data availability}

\textbf{We retained every reported participant-condition record that had the required source data and produced stride pairs under the common matching rules.} We did not remove records based on match count or stride-length error because both are study outcomes. Of the 128 possible records from 32 participants across four conditions, 126 entered the paired analysis. The GR Pulse source record was unavailable for ID\_186 L2 and ID\_195 L2, so those two conditions could not be aligned to the GaitRite timeline. Table~\ref{tab:s_excluded_pairs} lists these unavailable records.

The 126 available records contained 1{,}764 walks. We accepted a walk-level alignment when at least 60\% of the expected GaitRite heel strikes had an IMU match within 150~ms. A total of 1{,}669 walks passed this check and 95 were classified as synchronisation failures. These included 14 single-task, 13 serial-2, 19 serial-3, and 49 serial-7 walks. Failed alignments did not enter paired stride or event-agreement analyses. We did not classify them as sensor failures or remove any walk solely because of its stride-length error.

Three aligned walks did not enter stride pairing because a valid stride-level comparison was not possible. ID\_107 single-task trial 2 OUT and ID\_149 serial-2 trial 1 BACK contained an IMU event before mat entry that did not correspond to a GaitRite stride. ID\_195 serial-7 trial 1 BACK contained inconsistent GaitRite heel-strike and stride-time records. These walk-level decisions were defined in the shared analysis configuration and did not remove the remaining walks from those participant-condition records.

\begin{table}[!ht]
\centering
\caption{Reported-condition records unavailable for paired ZUPT-GaitRite analysis. L2 = serial subtraction by 2.}
\label{tab:s_excluded_pairs}
\footnotesize
\begin{tabular}{@{}llp{8.6cm}@{}}
\toprule
Subject & Condition & Diagnostic note \\
\midrule
ID\_186 & L2 & GR Pulse source record unavailable \\
ID\_195 & L2 & GR Pulse source record unavailable \\
\bottomrule
\end{tabular}
\end{table}

% ============================================================================
% S3: MoCA-based cognitive-status sensitivity reanalysis
% ============================================================================

\section*{S3. MoCA-based cognitive-status sensitivity reanalysis}

\textbf{We re-ran the cognitive-status stratification after excluding the one MoCA-defined CI participant whose score fell at the moderate-impairment boundary.} All 32 analysed participants had recorded MoCA scores (Table~\ref{tab:s_demographics}). The MoCA-defined CI group ($n = 13$) spanned MoCA 17--25. Twelve participants scored 18--25 and one scored 17, at the boundary of moderate impairment. Restricting the group to the twelve participants with MoCA scores of 18--25 returned an RMSE of 4.01~cm, compared with 4.14~cm in the full CI group (Table~\ref{tab:s_moca_sensitivity}). This result shows that the group estimate does not depend on the single moderate-range participant.

\begin{table}[!ht]
\centering
\caption{Cognitive-status sensitivity reanalysis. Primary = paper-reported MoCA-defined screening groups (all 32 participants, complete MoCA). Restricted CI = MoCA-defined CI group restricted to the twelve participants with MoCA 18--25, excluding the one participant at MoCA 17. MAE statistics are per-stride pooled.}
\label{tab:s_moca_sensitivity}
\small
\begin{tabular}{@{}llrrrrr@{}}
\toprule
Rule & Group & $n$ subj & $N$ pairs & RMSE (cm) & MAE (cm) & Bias (cm) \\
\midrule
Primary (paper)          & CN              & 19 & 6{,}139 & 3.65 & 2.29 & $+0.28$ \\
                         & CI              & 13 & 4{,}333 & 4.14 & 2.79 & $+0.24$ \\
\midrule
Restricted CI (sensitivity) & CN           & 19 & 6{,}139 & 3.65 & 2.29 & $+0.28$ \\
                         & CI (MoCA 18--25) & 12 & 3{,}963 & 4.01 & 2.71 & $+0.32$ \\
\bottomrule
\end{tabular}
\end{table}

% ============================================================================
% S4: Gravity-residual correction ablation
% ============================================================================

\section*{S4. Gravity-residual correction ablation}

\textbf{We re-ran the full pipeline on the same 32-participant cohort with the retroactive gravity-residual correction disabled.} The ablation kept stance detection, bias estimation, orientation tracking, and the per-stride zero-velocity reset unchanged. We removed only the swing-phase gravity-residual ramp. Both analyses used the same 10{,}472 paired strides. Disabling the correction increased pooled RMSE by 0.99~cm and shifted bias from $+0.3$ to $+2.2$~cm. The proportional-bias slope changed from $0.036$ to $0.044$, and ICC fell from 0.976 to 0.962 (Table~\ref{tab:s_ablation_smoothing}).

\begin{table}[!ht]
\centering
\caption{Gravity-residual correction ablation on the same 10{,}472 paired strides. ON = canonical pipeline and OFF = ablation. Bias = IMU minus GaitRite. Slope = proportional-bias regression of stride-level error on the mean of paired stride lengths.}
\label{tab:s_ablation_smoothing}
\small
\begin{tabular}{@{}lrrrrrr@{}}
\toprule
Condition & Correction & RMSE (cm) & MAE (cm) & Bias (cm) & ICC(2,1) & Slope \\
\midrule
\multirow{2}{*}{Pooled}      & OFF & 4.85 & 3.42 & $+2.25$ & 0.962 & 0.044 \\
                             & ON  & 3.86 & 2.50 & $+0.26$ & 0.976 & 0.036 \\
\midrule
\multirow{2}{*}{Single-task} & OFF & 5.25 & 3.98 & $+2.94$ & 0.955 & 0.052 \\
                             & ON  & 3.96 & 2.72 & $+0.71$ & 0.973 & 0.033 \\
\multirow{2}{*}{Serial 2s}   & OFF & 4.49 & 3.10 & $+1.84$ & 0.966 & 0.029 \\
                             & ON  & 3.70 & 2.39 & $+0.07$ & 0.977 & 0.033 \\
\multirow{2}{*}{Serial 3s}   & OFF & 4.78 & 3.23 & $+1.98$ & 0.963 & 0.024 \\
                             & ON  & 3.91 & 2.41 & $+0.03$ & 0.975 & 0.026 \\
\multirow{2}{*}{Serial 7s}   & OFF & 4.87 & 3.39 & $+2.26$ & 0.956 & 0.066 \\
                             & ON  & 3.87 & 2.49 & $+0.26$ & 0.971 & 0.048 \\
\bottomrule
\end{tabular}
\end{table}

% ============================================================================
% S5: Bootstrap robustness of cohort-level accuracy
% ============================================================================

\section*{S5. Bootstrap robustness of cohort-level accuracy}

\textbf{A participant-clustered bootstrap shows stable cohort-level accuracy under participant resampling.} We resampled the 32 analysed participants with replacement 10{,}000 times and recomputed each accuracy metric on the resulting concatenated stride pool. Because participants contributed different numbers of strides, the pool size varied across resamples. The 95\% percentile confidence intervals are shown in Table~\ref{tab:s_robustness}.

\begin{table}[!ht]
\centering
\caption{Cohort-level accuracy metrics with 95\% confidence intervals.}
\label{tab:s_robustness}
\small
\begin{tabular}{@{}lccc@{}}
\toprule
Metric & Point estimate & 95\% CI & CI half-width \\
\midrule
RMSE (cm)              & 3.86  & [3.51, 4.23]   & 0.36 \\
MAE (cm)               & 2.50  & [2.30, 2.70]   & 0.20 \\
Bias (cm)              & $+0.26$ & [$-0.08$, $+0.63$] & 0.35 \\
ICC(2,1)               & 0.976 & [0.964, 0.981] & 0.008 \\
\bottomrule
\end{tabular}
\end{table}

% ============================================================================
% S6: Stride-pair mismatch-filter sensitivity
% ============================================================================

\newpage
\needspace{12\baselineskip}
\section*{S6. Sensitivity to the stride-pair mismatch filter}

\textbf{Removing the stride-pair mismatch filter retained 22 additional strides and left typical absolute error similar.} This filter required IMU stride length and stride time to remain within 50--150\% of their matched GaitRite values. We retained all other participant, condition, walk, GaitRite plausibility, on-mat, and mutual-nearest matching rules. The resulting pool contained 10{,}494 rather than 10{,}472 strides. Pooled RMSE rose from 3.86 to 5.49~cm, while MAE changed from 2.50 to 2.65~cm (Table~\ref{tab:s_ratio_sensitivity}).

The filter removed 22 strides (0.21\% of the unfiltered pool). Their median stride-length error was $-72.2$~cm, RMSE was 85.2~cm, and maximum absolute error was 129.2~cm. These values are consistent with missed or merged stride events or unstable stride estimates rather than ordinary estimator error. The difference between pooled RMSE and MAE therefore reflects RMSE's sensitivity to a small number of large mismatches rather than a broad shift in typical stride accuracy.

\begin{table}[!ht]
\centering
\caption{Stride-length accuracy with and without the stride-pair mismatch filter. The filter restricted IMU stride length and stride time to 50--150\% of the matched GaitRite value. No mismatch filter = both checks omitted, with all other filtering and matching rules unchanged. Bias = IMU minus GaitRite.}
\label{tab:s_ratio_sensitivity}
\small
\begin{tabular}{@{}lrrrrr@{}}
\toprule
Analysis & $N$ pairs & RMSE (cm) & MAE (cm) & Bias (cm) & ICC(2,1) \\
\midrule
Mismatch filter    & 10{,}472 & 3.86 & 2.50 & $+0.26$ & 0.976 \\
No mismatch filter & 10{,}494 & 5.49 & 2.65 & $+0.17$ & 0.952 \\
\bottomrule
\end{tabular}
\end{table}

% ============================================================================
% S7: Gravity-residual correction gain sensitivity
% ============================================================================

\needspace{12\baselineskip}
\section*{S7. Gravity-residual correction gain sensitivity}

\textbf{Stride-length accuracy was stable across retroactive gravity-residual correction gains from 0.3 to 0.7.} We varied only the correction gain from 0.3 to 0.7 in increments of 0.1. All settings used the canonical stance thresholds, filtering, matching, and walk-alignment rules. RMSE changed smoothly from 4.10~cm at 0.3 to 3.83~cm at 0.6 and 3.87~cm at 0.7 (Table~\ref{tab:s_gain_sensitivity}). The canonical gain of 0.5 produced 3.86~cm RMSE, within 0.03~cm of the smallest value. The selected gain therefore lies in a broad near-optimal region rather than at a sharp optimum. Because all gains were evaluated in the same cohort, this analysis tests parameter stability but does not replace independent validation.

\begin{table}[!ht]
\centering
\caption{Sensitivity of pooled stride-length accuracy to the retroactive gravity-residual correction gain. Bias = IMU minus GaitRite.}
\label{tab:s_gain_sensitivity}
\small
\begin{tabular}{@{}lrrrrr@{}}
\toprule
Gain & $N$ pairs & RMSE (cm) & MAE (cm) & Bias (cm) & ICC(2,1) \\
\midrule
0.3 & 10{,}472 & 4.10 & 2.71 & $+1.05$ & 0.973 \\
0.4 & 10{,}472 & 3.95 & 2.57 & $+0.66$ & 0.975 \\
0.5 (canonical) & 10{,}472 & 3.86 & 2.50 & $+0.26$ & 0.976 \\
0.6 & 10{,}472 & 3.83 & 2.49 & $-0.13$ & 0.976 \\
0.7 & 10{,}472 & 3.87 & 2.54 & $-0.52$ & 0.976 \\
\bottomrule
\end{tabular}
\end{table}

\end{document}
